% Inserted after the scheduler section.
% Two additions: a bounded-budget selection rule, and phase exclusion.

\section{Selection under a bounded budget}
\label{sec:budget}

The scheduler of \Cref{sec:scheduler} runs every ready node until none
remains. That is the right default for a protocol submitted in full and
executed to completion, and it is what the guarantees of
\Cref{thm:scheduler} describe. It is the wrong default when the ready set is
larger than the resources available to it, which is the ordinary case once a
kernel is fed by many contributors at once.

This section adds a selection rule for that case. It is deliberately placed
\emph{outside} the kernel: the rule requires a per-node gain, gain is an
interpretive quantity, and by \Cref{thm:inertia} the kernel computes no such
thing. What follows is therefore a policy a module may apply when choosing
what to contribute, not a faculty the kernel acquires.

\begin{definition}[Budget, cost, gain]
\label{def:budget}
A \emph{budget} $\att>0$ bounds the total cost committed per unit interval:
if $R$ is the set of executions committed, $\sum_{r\in R}\cost(r)\le\att$.
Each candidate carries a cost $\cost(r)\ge\flo>0$, bounded below by the
floor, and a \emph{gain} $\gain(r)$ supplied by the submitting module. The
\emph{gain density} of a candidate is $\gain(r)/\cost(r)$.
\end{definition}

\begin{theorem}[Selection is a threshold on gain density]
\label{thm:threshold}
Order candidates by gain density and commit them in decreasing order until
the budget is exhausted. Equivalently, there is a price $p^{\star}\ge0$ such
that a candidate is committed iff its density is at least $p^{\star}$. Then:
\begin{enumerate}[label=\textup{(\roman*)},leftmargin=2.2em]
\item the rule is optimal for the continuous relaxation, and its value is
within the gain of a single candidate of the integer optimum;
\item if the repertoire is \emph{divisible} --- every candidate is an
independently selectable unit of cost $\flo$ --- the rule is exactly optimal
among integer selections.
\end{enumerate}
\end{theorem}

\begin{proof}
The problem is to maximise $\sum_{r\in R}\gain(r)$ subject to
$\sum_{r\in R}\cost(r)\le\att$. Taking candidates in decreasing density
until the budget is spent is optimal for the continuous relaxation, and the
Lagrangian dual assigns a single multiplier $p^{\star}$ such that a
candidate is selected iff its density exceeds it \citep{dantzig1957}. The
greedy solution differs from the continuous optimum only in the fractional
part of the boundary candidate, whose contribution is at most its own gain,
giving (i).

For (ii), if every candidate costs exactly $\flo$ then any budget admits
$\lfloor\att/\flo\rfloor$ candidates whatever they are, so the feasible set
is determined by cardinality alone; maximising the sum over a fixed
cardinality is achieved by taking the largest gains, which is what ordering
by density does when all costs are equal. The continuous optimum is
therefore attained at an integer point.
\end{proof}

\begin{remark}[Divisibility is the load-bearing hypothesis, and it is easy to
lose]
\label{rem:divisibility}
Clause (ii) requires candidates to be \emph{independently selectable} units,
not merely to have costs that are integer multiples of the floor. The
distinction is not pedantic: bundling several floor-sized units into one
indivisible candidate restores an ordinary $0/1$ knapsack, where the greedy
rule is only within one candidate of the optimum.

A witness, produced by the accompanying computation. Six candidates with
costs $(0.5,0.5,0.75,0.5,0.75,1.0)$ and gains
$(0.2,0.2,0.3,0.4,0.6,0.8)$ under a budget of $1.75$: the greedy rule
attains $1.2$ where the optimum is $1.4$. Every cost is a multiple of the
floor $0.25$, so the superficial reading of ``floor-additive'' is satisfied
and the exact claim nonetheless fails. We state clause (ii) for divisible
repertoires only.
\end{remark}

\begin{proposition}[The price rises as the budget narrows]
\label{prop:shadow-price}
Let $p^{\star}(\att)$ be the gain density of the highest-density candidate
the budget excludes. Then $p^{\star}$ is non-increasing in $\att$: widening
the budget weakly lowers the price, and narrowing it weakly raises it.
\end{proposition}

\begin{proof}
Candidates are considered in a fixed density order independent of $\att$.
Increasing the budget can only extend the prefix of that order which is
admitted, so the first excluded candidate moves later in the order or stays
where it is; since the order is decreasing in density, its density weakly
falls.
\end{proof}

\begin{remark}[Why the marginal excluded candidate, and not the last
admitted one]
\label{rem:price-definition}
The price must be defined as the density of the highest-density candidate
the budget \emph{excludes}. The density of the last candidate \emph{admitted}
is a different quantity and is not monotone: widening the budget can newly
admit a cheap, high-density candidate that did not previously fit, raising
that value while the true shadow price falls. The accompanying computation
exhibits the non-monotonicity directly. We record it because the two
definitions coincide often enough that the wrong one survives casual testing.
\end{remark}

\begin{corollary}[Declining is forced, not a defect]
\label{cor:declining}
Under a bounded budget a module must decline some candidates whenever the
total cost of the ready set exceeds $\att$. Since each cost is at least
$\flo>0$, at most $\lfloor\att/\flo\rfloor$ candidates can be committed per
interval regardless of their gains. Declining is therefore a consequence of
boundedness and not evidence of a scheduling failure.
\end{corollary}

\section{Phase exclusion}
\label{sec:phase}

A second bound on throughput is independent of the budget and follows from
the structure of contribution.

\begin{definition}[Construction and commitment]
\label{def:phases}
A module is \emph{constructing} when it is forming new subtask identities and
chunks to contribute, and \emph{committing} when it is executing chunks
already contributed.
\end{definition}

\begin{axiom}[Phase exclusion]
\label{ax:phase}
A module does not construct and commit in the same instant.
\end{axiom}

\begin{proposition}[Alternation, and a ceiling on commitment]
\label{prop:alternation}
Under \Cref{ax:phase}, a module that spends a fraction $\phi\in[0,1]$ of its
instants constructing commits in at most a fraction $1-\phi$ of them.
Consequently the committed record grows at most at rate $(1-\phi)$ per
instant, whatever the budget $\att$, and a module that never constructs
never acquires new nodes to commit to.
\end{proposition}

\begin{proof}
By \Cref{ax:phase} the two phases partition the instants a module occupies,
so the fraction available for commitment is $1-\phi$. The committed record
increments only on commitment (\Cref{def:record}), giving the rate bound.
The final clause is immediate: nodes enter the graph only by contribution,
which occurs in the construction phase.
\end{proof}

\begin{remark}[Two independent reasons a sweep quiesces]
\label{rem:two-ceilings}
\Cref{cor:declining} and \Cref{prop:alternation} are separate. The first is
economic: the budget cannot pay for every ready candidate. The second is
structural: an interval spent constructing is not available for committing,
and no budget relaxes it. A kernel that observes low throughput cannot
distinguish the two from the record alone, and a module reporting one should
not be read as reporting the other.

Neither ceiling gives the kernel a verdict. A declined candidate is not a
failed one, and an instant spent constructing is not an idle one; both are
recorded as what they are, and \Cref{thm:inertia} is untouched.
\end{remark}
